%% ch03-first-steps.tex — JA4proxy Operator's User Guide
%% Chapter 3: First Steps

\chapter{First Steps}
\label{ch:firststeps}
\index{first steps}

\newthought{The proxy is running and the dashboard is open.} Now what?

This chapter walks through what you will see in the first minutes of operation:
what the logs mean, how to generate representative test traffic, and how to
read the Grafana dashboard. By the end of this chapter you should be able to
tell from the dashboard whether \ja\ is classifying traffic correctly.

\section{What Happens at Dial Zero}
\index{dial!monitor mode}
\index{monitor mode}

At \code{dial=0} (the default), \ja\ is in monitor mode. Every connection is
analysed, scored, and logged — but no connection is ever blocked. The proxy
passes everything through to your backend regardless of its risk score.

This is the right mode for initial deployment. You are building a baseline
understanding of what your traffic looks like before you start acting on it.

Logs will show entries like this:

\begin{lstlisting}[language=bash, numbers=none]
ALLOWED:  172.19.0.10 | Country: IE | JA4: t13d1113h2_8daaf6152771_b0da82dd1658 |
          Name: Browser (TLS 1.3) | TLS: TLS 1.3
ALLOWED:  185.220.0.1 | Country: RU | JA4: t13d0912_a1b2c3d4e5f6_987654321abc |
          Name: Tool/Bot (TLS 1.3) | TLS: TLS 1.3 | Score: 74
\end{lstlisting}

At dial=0, both lines show \code{ALLOWED}. The second entry has a high risk score (74)
and would be blocked at dial=100 — but in monitor mode it passes through. This is
correct behaviour. You are observing, not acting.\sidenote{The score is still
computed and recorded even at dial=0. This means when you eventually raise the dial,
you already have weeks of score history to reason about.}

\section{The Traffic Generator}
\index{traffic generator}
\index{generate-tls-traffic.sh}

\ja\ ships with a traffic generator for testing and demonstration. It sends realistic
TLS connections using a mix of legitimate browser profiles and known-malicious tool
profiles.\sidenote{Profiles included: Chrome, Firefox, Safari (legitimate) and Sliver
C2, Cobalt Strike, Evilginx, Python bot, credential stuffer (malicious).}

\begin{bashcode}
./scripts/generate-tls-traffic.sh <duration_seconds> <legit_percent> <workers>
\end{bashcode}

For example, to run for 60 seconds with 30\% legitimate traffic and 10 workers:

\begin{bashcode}
./scripts/generate-tls-traffic.sh 60 30 10
\end{bashcode}

A good first run to populate the dashboard:

\begin{bashcode}
# 60 seconds, mixed traffic, moderate load
./scripts/generate-tls-traffic.sh 60 50 20
\end{bashcode}

While it runs you will see log output showing connections being processed. Some will
have fingerprints identified as known tools; these will show elevated risk scores.
With dial=0 they will all be allowed, but the scoring information appears in the log.

\section{Reading the Logs}
\index{log format}
\index{logs!reading}

\subsection{Log Line Format}

Each connection produces one log line on close. The format is:

\begin{lstlisting}[language=bash, numbers=none]
{ACTION}: {client_ip} | Country: {cc} | JA4: {fingerprint} |
          Name: {label} | TLS: {version} [| Reason: {reason}]
          [| Score: {score}]
\end{lstlisting}

\begin{table}[ht]
  \small
  \begin{tabularx}{\linewidth}{lX}
    \toprule
    \textbf{Field} & \textbf{Meaning} \\
    \midrule
    \code{ACTION} & \code{ALLOWED}, \code{BLOCKED}, \code{TARPITTED}, or \code{RATE\_LIMITED} \\
    \code{client\_ip} & Real client IP (extracted from PROXY protocol header) \\
    \code{Country} & Two-letter country code from GeoIP lookup \\
    \code{JA4} & The computed JA4 fingerprint string \\
    \code{Name} & Human-readable label for the fingerprint, if known \\
    \code{TLS} & TLS version negotiated in the ClientHello \\
    \code{Reason} & For blocked connections, the blocking reason \\
    \code{Score} & Risk score (0--100); appears when score is non-zero \\
    \bottomrule
  \end{tabularx}
  \caption{Log line fields}
\end{table}

\subsection{Example Log Lines}

A legitimate browser connecting from Ireland:
\begin{lstlisting}[language=bash, numbers=none]
ALLOWED:  172.19.0.10 | Country: IE | JA4: t13d1113h2_8daaf6152771_b0da82dd1658 |
          Name: Browser (TLS 1.3) | TLS: TLS 1.3
\end{lstlisting}

A known tool blocked by the JA4 blacklist:
\begin{lstlisting}[language=bash, numbers=none]
BLOCKED:  185.220.0.1 | Country: RU | JA4: t13d0912_a1b2c3d4e5f6_987654321abc |
          Name: Tool/Bot (TLS 1.3) | Reason: JA4 blacklisted
\end{lstlisting}

An IP that has been banned (reached rate limit and been automatically banned):
\begin{lstlisting}[language=bash, numbers=none]
BLOCKED:  10.0.0.42 | Country: DE | JA4: t13d1113h2_8daaf6152771_b0da82dd1658 |
          Name: Browser (TLS 1.3) | Reason: Banned for 300s
\end{lstlisting}

A high-scoring connection at dial=0 (passed through but scored):
\begin{lstlisting}[language=bash, numbers=none]
ALLOWED:  198.51.100.3 | Country: NL | JA4: t13d0809_c0ffee123456_aabbccddeeff |
          Name: Python-requests (TLS 1.3) | TLS: TLS 1.3 | Score: 68
\end{lstlisting}

\subsection{Accessing Logs}

Live tail all proxy output:
\begin{bashcode}
docker compose logs -f ja4proxy
\end{bashcode}

Filter for blocked connections only:
\begin{bashcode}
docker compose logs ja4proxy | grep "^BLOCKED"
\end{bashcode}

Search for a specific IP:
\begin{bashcode}
docker compose logs ja4proxy | grep "185.220.0.1"
\end{bashcode}

\section{The Grafana Dashboard}
\index{Grafana}
\index{dashboard}

Open \url{http://localhost:3001} and navigate to the \emph{JA4proxy Overview} dashboard.
After running the traffic generator you should see all panels populated.

\subsection{Panel-by-Panel Guide}

\subsection{Connection Rate}
\index{Grafana!connection rate panel}

Shows connections per second over time, coloured by action:
\begin{itemize}
  \item Green: \badgeallow{} — connections passed to backend
  \item Red: \badgeblock{} — connections rejected
  \item Amber: \badgetarpit{} — connections held in tarpit
  \item Dark: \badgeban{} — connections from currently-banned IPs
\end{itemize}

At dial=0, everything is green. When you raise the dial, you will start to see
red and amber segments appear.

\subsection{Action Distribution}
\index{Grafana!action distribution}

A pie or bar chart showing the percentage breakdown of actions over the selected
time window. This is the key panel for understanding how \ja\ is classifying
your traffic. In a healthy deployment you should see the vast majority as ALLOWED,
with a small but non-zero fraction as BLOCKED.

If you see zero BLOCKED and you know malicious traffic is present, the dial is
probably at zero. If you see a high BLOCKED fraction, investigate before raising
the dial further — high block rates often mean a fingerprint list needs tuning.

\subsection{Risk Score Distribution}
\index{Grafana!risk score distribution}

A histogram of risk scores across all connections. This is your most important
panel for dial calibration. The shape of this histogram tells you where to set
blocking thresholds:

\begin{itemize}
  \item A bimodal distribution (peaks near 0 and near 80--100) is ideal: legitimate
    traffic clusters at low scores, bad traffic clusters at high scores.
  \item A single peak near zero means either the traffic is mostly legitimate, or
    the signal collectors are not seeing the bad traffic you expect.
  \item A broad flat distribution suggests mixed signals — worth investigating
    which specific signals are contributing.
\end{itemize}

\subsection{Top JA4 Fingerprints}
\index{Grafana!fingerprint panel}

A table of the most frequent JA4 fingerprints seen, with their labels (if known),
connection counts, and average risk scores. This is where you identify candidate
fingerprints for your whitelist or blacklist.

If a fingerprint appears frequently with a high average score but a human-readable
label that sounds legitimate, investigate before blacklisting it — it may be a
shared fingerprint used by both automated and manual clients.

\subsection{Geographic Distribution}

A world map or table showing connection volume by country. Useful for identifying
whether a sudden traffic increase is concentrated in an unexpected region.

\subsection{Rate Limiting Activity}

Shows how many connections are hitting each rate-limit strategy (by IP, by JA4
fingerprint, by IP+JA4 pair). High activity here means automated clients are
testing your rate limits.

\subsection{Redis Lag}

Shows Redis response time. Spikes here affect scoring latency. If Redis lag is
consistently above 1ms, investigate the Redis container's resource usage.

\section{Understanding Fingerprint Labels}
\index{fingerprint labels}
\index{JA4!labels}

The \emph{Name} field in logs and the label column in the Grafana fingerprint table
come from \code{config/proxy.yml}'s \configkey{fingerprint\_labels} section. This
is a mapping from JA4 fingerprint strings to human-readable names.

\ja\ ships with labels for common browsers, known security tools, and common C2
frameworks. You will extend this list as you observe your own traffic.

When you see an unfamiliar fingerprint appearing frequently, look it up in the JA4
community database at \url{https://ja4db.com} and add an appropriate label to your
config. This makes the dashboard and logs much easier to interpret.

\section{Your First Configuration Change}
\index{configuration!first change}

The two most common initial configuration tasks are:

\subsection{Setting Your Backend Host}

If you have not already set \code{BACKEND\_HOST} in \code{.env}, the proxy is
forwarding to the default test backend. Update \code{.env} and restart:

\begin{bashcode}
# In .env:
BACKEND_HOST=prod.internal.example.com

make stop
make start
\end{bashcode}

\subsection{Enabling Country Filtering}

To restrict traffic to specific countries, edit \code{config/proxy.yml}:

\begin{yamlcode}
# Enable the whitelist (only listed countries pass)
country_whitelist_enabled: true
country_whitelist:
  - IE
  - GB
  - US
  - DE
\end{yamlcode}

Then reload the config without restarting (see Section~\ref{sec:hotreload}):

\begin{bashcode}
# Signal the proxy to reload config
kill -HUP $(docker compose exec ja4proxy cat /var/run/proxy.pid)
\end{bashcode}

\begin{warningbox}[Country filtering and false positives]
Country whitelisting is a blunt instrument. VPN users, Tor exit nodes, and CDN
origin IPs may appear to come from unexpected countries. Enable country filtering
only if you have a genuine operational reason to restrict by geography, and verify
your legitimate user base is not affected before enabling.
\end{warningbox}
